\begin{frame}[standout]

    Avez-vous été attentifs ?

    À vous de jouer !

\end{frame}

{
\setbeamercovered{transparent}

\begin{frame}{Q1}

    \textbf{Qu'est-ce que le NoSQL ?}

    \begin{itemize}
        \item \uncover<1>{\textcolor{kh-red}{Langage de programmation de manipulation de bases de données relationnelles}}
        \item \textbf<2>{\textcolor{kh-blue}{Bases de données non relationnelles conçues pour des données distribuées et volumineuses}}
        \item \uncover<1>{\textcolor{kh-green}{Protocole de communication pour les réseaux informatiques}}
        \item \uncover<1>{\textcolor{kh-yellow}{Système de gestion de bases de données exclusivement destiné au traitement des transactions}}
    \end{itemize}

\end{frame}

\begin{frame}{Q2}

    \textbf{Quelle est la principale utilité du partitionnement dans les bases de données NoSQL ?}

    \begin{itemize}
        \item \textbf<2>{\textcolor{kh-red}{Il améliore les performances en répartissant les données sur plusieurs nœuds}}
        \item \uncover<1>{\textcolor{kh-blue}{Il sert uniquement à réduire la taille physique des fichiers de la base de données sur le disque dur}}
        \item \uncover<1>{\textcolor{kh-green}{Il permet de crypter automatiquement les données pour renforcer la sécurité des informations stockées}}
        \item \uncover<1>{\textcolor{kh-yellow}{Il est utilisé pour convertir les bases de données NoSQL en bases de données relationnelles afin de supporter le langage SQL}}
    \end{itemize}

\end{frame}

\begin{frame}{Q2}

    \textbf{Quelle est la principale caractéristique des bases de données NoSQL ?}

    \begin{itemize}
        \item \textbf<2>{\textcolor{kh-red}{Propriétés BASE}}
        \item \uncover<1>{\textcolor{kh-blue}{Transactions ACID}}
        \item \uncover<1>{\textcolor{kh-green}{Langage SQL}}
        \item \uncover<1>{\textcolor{kh-yellow}{Open Database Connectivity}}
    \end{itemize}

\end{frame}

\begin{frame}{Q3}

    \textbf{Quels sont les mécanismes pour assurer disponibilité et tolérance aux pannes dans un environnement distribué ?}

    \begin{itemize}
        \item \textbf<2>{\textcolor{kh-red}{Réplication des données}}
        \item \textbf<2>{\textcolor{kh-blue}{Sharding}}
        \item \uncover<1>{\textcolor{kh-green}{Transactions ACID strictes}}
        \item \textbf<2>{\textcolor{kh-yellow}{Cohérence éventuelle}}
    \end{itemize}

\end{frame}

\begin{frame}{Q4}

    \textbf{Que signifient les initiales du théorème CAP ?}

    \begin{itemize}
        \item \uncover<1>{\textcolor{kh-red}{Complexity, Accuracy, Performance}}
        \item \uncover<1>{\textcolor{kh-blue}{Centralized Access Protocol}}
        \item \textbf<2>{\textcolor{kh-green}{Consistency, Availability, Partition Tolerance}}
        \item \uncover<1>{\textcolor{kh-yellow}{Consistency, Availability, Performance}}
    \end{itemize}

\end{frame}

%\begin{frame}
%
%    Ajouter des questions ouvertes (ex : "Proposez une base NoSQL adaptée à un système de recommandation de films").
%    Utiliser des cas réels (ex : "Quel type de base utiliseriez-vous pour stocker les logs d’un site web ?").
%
%\end{frame}

%\begin{frame}{Q5}
%
%    \textbf{}
%
%    \begin{itemize}
%        \item \uncover<1>{\textcolor{kh-red}{}}
%        \item \uncover<1>{\textcolor{kh-blue}{}}
%        \item \textbf<2>{\textcolor{kh-green}{}}
%        \item \uncover<1>{\textcolor{kh-yellow}{}}
%    \end{itemize}
%
%\end{frame}
%
%\begin{frame}{Q6}
%
%    \textbf{}
%
%    \begin{itemize}
%        \item \uncover<1>{\textcolor{kh-red}{}}
%        \item \uncover<1>{\textcolor{kh-blue}{}}
%        \item \textbf<2>{\textcolor{kh-green}{}}
%        \item \uncover<1>{\textcolor{kh-yellow}{}}
%    \end{itemize}
%
%\end{frame}
%
%\begin{frame}{Q7}
%
%    \textbf{}
%
%    \begin{itemize}
%        \item \uncover<1>{\textcolor{kh-red}{}}
%        \item \uncover<1>{\textcolor{kh-blue}{}}
%        \item \textbf<2>{\textcolor{kh-green}{}}
%        \item \uncover<1>{\textcolor{kh-yellow}{}}
%    \end{itemize}
%
%\end{frame}
%
%\begin{frame}{Q8}
%
%    \textbf{}
%
%    \begin{itemize}
%        \item \uncover<1>{\textcolor{kh-red}{}}
%        \item \uncover<1>{\textcolor{kh-blue}{}}
%        \item \textbf<2>{\textcolor{kh-green}{}}
%        \item \uncover<1>{\textcolor{kh-yellow}{}}
%    \end{itemize}
%
%\end{frame}
%
%\begin{frame}{Q9}
%
%    \textbf{}
%
%    \begin{itemize}
%        \item \uncover<1>{\textcolor{kh-red}{}}
%        \item \uncover<1>{\textcolor{kh-blue}{}}
%        \item \textbf<2>{\textcolor{kh-green}{}}
%        \item \uncover<1>{\textcolor{kh-yellow}{}}
%    \end{itemize}
%
%\end{frame}
%
%\begin{frame}{Q10}
%
%    \textbf{}
%
%    \begin{itemize}
%        \item \uncover<1>{\textcolor{kh-red}{}}
%        \item \uncover<1>{\textcolor{kh-blue}{}}
%        \item \textbf<2>{\textcolor{kh-green}{}}
%        \item \uncover<1>{\textcolor{kh-yellow}{}}
%    \end{itemize}
%
%\end{frame}

}
